\documentclass[11pt,a4paper]{article}

% ---------- fontes e língua ----------
\usepackage[T1]{fontenc}
\usepackage[utf8]{inputenc}
\usepackage{lmodern}
\usepackage{tgpagella}
\usepackage{tgheros}
\usepackage[portuguese]{babel}
\usepackage{microtype}

% ---------- página e cores ----------
\usepackage[a4paper,top=2.6cm,bottom=2.4cm,left=2.4cm,right=2.4cm]{geometry}
\usepackage[table]{xcolor}
\definecolor{azul}{RGB}{28,62,112}
\definecolor{azulclaro}{RGB}{226,234,246}
\definecolor{cinza}{RGB}{95,102,115}
\definecolor{laranja}{RGB}{190,105,25}
\definecolor{laranjaclaro}{RGB}{253,241,226}

% ---------- tabelas, listas, caixas ----------
\usepackage{booktabs}
\usepackage{array}
\usepackage{tabularx}
\usepackage{enumitem}
\usepackage{amssymb}
\usepackage{titlesec}
\usepackage{fancyhdr}
\usepackage{lastpage}
\usepackage{tcolorbox}

\setlength{\parindent}{0pt}
\setlength{\parskip}{6pt}

\setlist{nosep, leftmargin=*, itemsep=3pt}
\setlist[itemize]{label=\textcolor{azul}{\textbullet}}
\setlist[enumerate]{label=\textcolor{azul}{\sffamily\bfseries\arabic*.}}

\titleformat{\section}{\sffamily\Large\bfseries\color{azul}}{\thesection}{0.7em}{}[{\color{azul}\titlerule[0.6pt]}]
\titleformat{\subsection}{\sffamily\large\bfseries\color{azul}}{\thesubsection}{0.6em}{}
\titlespacing*{\section}{0pt}{1.6em}{0.8em}
\titlespacing*{\subsection}{0pt}{1.1em}{0.3em}

\pagestyle{fancy}
\setlength{\headheight}{15pt}
\fancyhf{}
\renewcommand{\headrule}{{\color{azul}\hrule width\headwidth height\headrulewidth}}
\renewcommand{\headrulewidth}{0.4pt}
\fancyhead[L]{\sffamily\small\color{cinza} Universidade da Beira Interior}
\fancyhead[R]{\sffamily\small\color{cinza} Engenharia de Software \textperiodcentered{} 2026/27}
\fancyfoot[L]{\sffamily\small\color{cinza} Guia de Jira}
\fancyfoot[R]{\sffamily\small\color{cinza} Página \thepage\ de \pageref{LastPage}}

\fancypagestyle{primeira}{%
  \fancyhf{}%
  \renewcommand{\headrulewidth}{0pt}%
  \fancyfoot[L]{\sffamily\small\color{cinza} Guia de Jira}%
  \fancyfoot[R]{\sffamily\small\color{cinza} Página \thepage\ de \pageref{LastPage}}%
}

\newcolumntype{Y}{>{\raggedright\arraybackslash}X}
\newcolumntype{L}[1]{>{\raggedright\arraybackslash}p{#1}}
\newcolumntype{N}[1]{>{\raggedright\arraybackslash\sffamily\bfseries}p{#1}}
\newcommand{\th}[1]{\textcolor{white}{\sffamily\bfseries #1}}

% nome de botão ou menu do Jira
\newcommand{\btn}[1]{{\setlength{\fboxsep}{1.6pt}\colorbox{azulclaro}{\sffamily\small\strut #1}}}
\newcommand{\seta}{\,$\rightarrow$\,}

% tabelas com cabeçalho azul e linhas alternadas
\newcommand{\tabelaini}{\begingroup\small\renewcommand{\arraystretch}{1.45}\rowcolors{2}{azulclaro!55}{white}}
\newcommand{\tabelafim}{\endgroup}

\tcbset{
  caixa/.style={colback=white, colframe=azul, coltitle=white, colbacktitle=azul,
    fonttitle=\sffamily\bfseries, arc=3pt, boxrule=0.6pt,
    left=9pt, right=9pt, top=6pt, bottom=6pt},
  dica/.style={colback=azulclaro!60, colframe=azul!70, coltitle=white, colbacktitle=azul!80,
    fonttitle=\sffamily\bfseries\small, arc=3pt, boxrule=0.5pt,
    left=9pt, right=9pt, top=5pt, bottom=5pt},
  atencao/.style={colback=laranjaclaro, colframe=laranja, coltitle=white, colbacktitle=laranja,
    fonttitle=\sffamily\bfseries\small, arc=3pt, boxrule=0.5pt,
    left=9pt, right=9pt, top=5pt, bottom=5pt}
}
\newtcolorbox{codigo}{colback=black!4, colframe=black!25, boxrule=0.4pt, arc=2pt,
  left=8pt, right=8pt, top=4pt, bottom=4pt, fontupper=\ttfamily\small}

\begin{document}
\thispagestyle{primeira}

% ---------- cabeçalho do documento ----------
\begin{tcolorbox}[colback=azul, colframe=azul, coltext=white, arc=3pt, boxrule=0pt,
    left=16pt, right=16pt, top=14pt, bottom=14pt]
  \sffamily
  {\footnotesize\bfseries\MakeUppercase{Universidade da Beira Interior}}\\[2pt]
  {\footnotesize Faculdade de Engenharia \textperiodcentered{} Departamento de Informática}\\[16pt]
  {\fontsize{28}{32}\selectfont\bfseries Guia de Jira}\\[4pt]
  {\large Como começar a usar o nosso Jira \textperiodcentered{} Engenharia de Software 2026/27}
\end{tcolorbox}

O Jira da equipa já está montado. Este guia é para quem vai usá-lo todos os dias e nunca o tinha visto: como entrar, onde está cada coisa, o que fazer em cada situação e o que tem de ficar escrito lá para o Scrum contar na avaliação. Não precisas de ler tudo. Faz primeiro os «primeiros 15 minutos» e volta aqui sempre que tiveres uma dúvida.

\begin{tcolorbox}[caixa, title={O essencial em poucas linhas}, colback=azulclaro!60]
  \begin{itemize}
    \item Cada coisa que há para fazer é um \textbf{cartão}. Se não está num cartão, para o Jira não existe.
    \item O trabalho do sprint está no \textbf{Quadro}, em três colunas: \textit{A fazer}, \textit{Fazendo} e \textit{Feito}.
    \item Pegas num cartão, pões lá o teu nome, arrastas para \textit{Fazendo} e, quando estiver mesmo pronto, para \textit{Feito}.
    \item Escreves a tua daily num comentário, todos os dias úteis.
    \item O que decidimos fica escrito no Jira. O que fica só no WhatsApp perde-se.
    \item Dúvidas de acesso ou de funcionamento: fala com o Tiago, que é o administrador do Jira.
  \end{itemize}
\end{tcolorbox}

%--------------------------------------------------------------------
\section{Os primeiros 15 minutos}

\subsection{Entrar no Jira}

\begin{enumerate}
  \item Procura na tua caixa de email uma mensagem da Atlassian com o convite do Tiago para o Jira. Se não a vires, olha também para o spam.
  \item Abre-a e carrega no botão para aceitar o convite. Se ainda não tens conta Atlassian, o Jira pede-te para a criares: usa o teu email e escolhe uma palavra-passe só tua, que não partilhas com ninguém.
  \item Depois de entrares, abre o endereço do nosso Jira: \texttt{ubi-team-ebw5rr6y.atlassian.net}.
  \item No menu \btn{Projetos}, escolhe \textit{Engenharia de Software}. A chave do projeto é \texttt{EDS}, por isso os cartões chamam-se \texttt{EDS-1}, \texttt{EDS-2}, e assim por diante.
  \item Guarda o endereço nos favoritos do browser, ou marca o projeto com a estrela, para o encontrares sempre à mão.
\end{enumerate}

\begin{tcolorbox}[atencao, title={Não recebeste o convite?}]
  Diz ao Tiago o email com que queres entrar. Ele convida-te de novo em \btn{Acesso}, nas definições do projeto. Quando o convite chega, tem de ser aceite com esse mesmo email.
\end{tcolorbox}

\subsection{Acertar o perfil}

\begin{itemize}
  \item \textbf{O teu nome.} O Jira mostra o nome que está na tua conta Atlassian, não o do contrato. Para toda a gente te reconhecer quando atribui cartões, muda-o para o teu nome real: clica na tua fotografia (ou nas iniciais), no canto superior direito, e procura a gestão da conta e do perfil.
  \item \textbf{A língua.} O nosso projeto está em português. Se preferires ver o Jira em inglês, nas definições pessoais da conta há uma opção de idioma. Os botões mudam de nome, mas fazem o mesmo (há uma tabela com os dois no ponto 2.2).
\end{itemize}

\subsection{Três coisas para fazeres já}

São tarefas a sério, não exercícios. Ajudam-te a perceber o Jira e adiantam trabalho da equipa.

\begin{itemize}[label=\textcolor{azul}{$\square$}]
  \item Abre o \btn{Quadro} e vê as três colunas. Procura os cartões que já têm o teu nome (se houver).
  \item Abre o cartão \textbf{EDS-24, Membros da equipa}, clica na descrição para a editar e escreve o teu número de aluno ao lado do teu nome.
  \item Abre o cartão \textbf{EDS-16, Certificado Scrum Foundation Professional}. Anexa lá o teu certificado, ou envia-o ao Tiago, para ir para a pasta \texttt{certificados-scrum} do repositório.
  \item Lê o contrato de equipa: está anexado ao cartão \textbf{EDS-1, Team Contract}, e o resumo está na página \textit{Team Contract} do Confluence.
\end{itemize}

%--------------------------------------------------------------------
\section{Conhecer o nosso Jira}

\subsection{A barra lateral}

Quando entras no projeto, a barra do lado esquerdo tem os sítios principais. Os nomes podem variar um pouco de uma versão para outra.

\tabelaini
\begin{tabularx}{\textwidth}{N{3cm} Y}
  \rowcolor{azul} \th{Sítio} & \th{Para que serve} \\
  Resumo & Visão geral do projeto e da atividade recente. \\
  Linha do tempo & Os epics numa linha do tempo, tipo roadmap. \\
  Backlog & A lista de tudo o que falta fazer e o planeamento dos sprints. É o ecrã do Product Owner. \\
  Quadro & O sprint em curso, com as colunas. É onde passamos o dia a dia. \\
  Relatórios & Gráficos do sprint (por exemplo o \textit{burndown}). Servem de prova na Review. \\
  Definições & Acessos e funcionalidades do projeto. Só os administradores mexem aqui. \\
  \bottomrule
\end{tabularx}
\tabelafim

\begin{tcolorbox}[dica, title={Dica: atalhos de teclado}]
  Dentro do Jira, a tecla \texttt{c} costuma abrir a criação de um cartão e a tecla \texttt{/} abre a pesquisa. A tecla \texttt{[} esconde e mostra a barra lateral.
\end{tcolorbox}

\subsection{As palavras do Jira}

\tabelaini
\begin{tabularx}{\textwidth}{N{2.8cm} L{3.4cm} Y}
  \rowcolor{azul} \th{Neste guia} & \th{Em inglês} & \th{O que é} \\
  Cartão & \textit{Issue}, \textit{Work item} & Uma coisa para fazer ou para registar. \\
  Tipo de cartão & \textit{Issue type} & \textit{Tarefa}, \textit{História}, \textit{Bug}, \textit{Subtask} ou \textit{Epic}. \\
  Epic & \textit{Epic} & Um cartão grande que agrupa outros, como uma pasta. \\
  Backlog & \textit{Backlog} & A lista do que falta fazer, por ordem de prioridade. \\
  Sprint & \textit{Sprint} & Um período fechado, com objetivo, em que fazemos parte do backlog. \\
  Quadro & \textit{Board} & As colunas do sprint em curso. \\
  Responsável & \textit{Assignee} & A pessoa que vai fazer o cartão. \\
  Estimativa & \textit{Story point estimate} & Um número que mede o esforço do cartão. \\
  Etiqueta & \textit{Label} & Uma marca livre para filtrar cartões. \\
  Sinalizar & \textit{Flag} & Destacar um cartão que está bloqueado. \\
  \bottomrule
\end{tabularx}
\tabelafim

\subsection{Como é um cartão}

Quando abres um cartão, vês estes campos. Alguns preenches tu, outros preenche-os o Jira ou o Planning.

\tabelaini
\begin{tabularx}{\textwidth}{N{3.3cm} Y L{3.3cm}}
  \rowcolor{azul} \th{Campo} & \th{O que é} & \th{Quem preenche} \\
  Chave (ex.: EDS-13) & O código único do cartão. Usamo-lo no GitHub. & O Jira \\
  Resumo & O título, curto e claro. & Quem cria \\
  Descrição & O que é para fazer, a user story e os critérios de aceitação. & Quem cria, com o PO \\
  Estado & Em que coluna está: \textit{A fazer}, \textit{Fazendo} ou \textit{Feito}. & Quem trabalha nele \\
  Responsável & Quem vai fazer o cartão. Uma só pessoa. & No Planning, ou tu ao pegares nele \\
  Story point estimate & A estimativa do esforço. & A equipa, em conjunto \\
  Principal (\textit{Parent}) & O epic a que pertence. & Quem cria \\
  Sprint & O sprint em que está. & No Planning \\
  Etiquetas & Marcas livres, como \textit{documento} ou \textit{scrum}. & Quem cria \\
  Comentários & A conversa sobre o cartão. & Todos \\
  Anexos & Ficheiros anexados: capturas, PDF, imagens. & Todos \\
  \bottomrule
\end{tabularx}
\tabelafim

\subsection{Como o nosso projeto está organizado}

O enunciado pede um Jira com várias secções. No nosso, algumas são ecrãs do próprio Jira e outras são epics com cor própria:

\tabelaini
\begin{tabularx}{\textwidth}{L{3.6cm} Y}
  \rowcolor{azul} \th{O enunciado pede} & \th{Onde está} \\
  Product Backlog & O ecrã \btn{Backlog}. \\
  To Do, In Progress, Done & As colunas do \btn{Quadro}: \textit{A fazer}, \textit{Fazendo}, \textit{Feito}. \\
  Team Contract & Epic \textbf{EDS-1} (roxo). Tem o contrato e o registo de uso de IA. \\
  Artifacts & Epic \textbf{EDS-2} (azul-petróleo). Links e cartões dos wireframes, mockups, arquitetura, modelação e repositório. \\
  Daily SCRUM & Epic \textbf{EDS-3} (verde). Os registos de Daily, Planning, Review e Retrospetiva. \\
  Presenças e faltas & Epic \textbf{EDS-17} (laranja). Um cartão por sprint. \\
  Equipa e papéis & Epic \textbf{EDS-18} (azul). Membros, rotação de PO e SM e certificados. \\
  Áreas do produto & Epics \textbf{EDS-48 a EDS-54} (vermelho), copiados do enunciado. \\
  \bottomrule
\end{tabularx}
\tabelafim

Em cada sprint há ainda cinco cartões de registo: \textit{Daily}, \textit{Planning}, \textit{Review}, \textit{Retrospetiva} e \textit{Presenças}. Estão dentro do sprint a que pertencem e passam a \textit{Feito} quando a reunião acontece e fica registada.

Os cartões têm etiquetas para filtrar: \textit{scrum}, \textit{presencas}, \textit{contrato}, \textit{ia}, \textit{documento}, \textit{equipa}, \textit{produto} e \textit{enunciado}.

\begin{tcolorbox}[dica, title={Também há o Confluence}]
  O Jira criou um espaço no Confluence com o nome \textit{Engenharia de Software}. Tem uma página inicial, o resumo do contrato, a equipa e os papéis, o SCRUM e os artefactos. Serve para leres as coisas com calma. O trabalho do dia a dia faz-se no Jira.
\end{tcolorbox}

\subsection{As colunas e as regras do quadro}

\begin{itemize}
  \item \textbf{A fazer:} está no sprint, mas ainda ninguém pegou nele.
  \item \textbf{Fazendo:} alguém está a trabalhar nele. Fica aqui até o \textit{pull request} ser aprovado e integrado.
  \item \textbf{Feito:} cumpre a nossa Definição de Concluído e o Product Owner aceitou.
\end{itemize}

O contrato tem duas regras para o quadro: ninguém tem mais de dois cartões abertos ao mesmo tempo, e um cartão só entra em \textit{Fazendo} depois de ter responsável e estimativa. No início ainda estamos a estimar, por isso usamos o bom senso e perguntamos no grupo.

%--------------------------------------------------------------------
\section{O dia a dia, passo a passo}

\subsection{Encontrar os teus cartões}

No \btn{Quadro}, em cima, vês as fotografias (ou iniciais) das pessoas da equipa. Carrega na tua e o quadro mostra só os cartões que são teus. Se quiseres uma lista, abre a pesquisa avançada (\btn{Filtros}) e cola:

\begin{codigo}
project = EDS AND assignee = currentUser() AND statusCategory != Done
\end{codigo}

\subsection{Pegar num cartão}

\begin{enumerate}
  \item Abre o cartão e carrega em \btn{Atribuir a mim}, ao lado de \textit{Responsável}.
  \item Arrasta-o no quadro de \textit{A fazer} para \textit{Fazendo}. Também podes mudar o estado no botão que está no topo do cartão.
  \item Se ainda não tem estimativa, diz na conversa do grupo, para a equipa estimar.
\end{enumerate}

\subsection{Atualizar enquanto trabalhas}

\begin{itemize}
  \item \textbf{Comenta} o que foste fazendo, mesmo que seja pouco. Quem abrir o cartão percebe logo onde estás.
  \item \textbf{Anexa} capturas de ecrã, ficheiros e links. Quanto mais claro, mais fácil é o Product Owner aceitar na Review.
  \item \textbf{Menciona} colegas com \texttt{@nome} para os chamar à conversa. Eles recebem um aviso.
  \item Se o trabalho tem partes distintas, cria \textbf{subtarefas} e atribui cada uma a quem a vai fazer.
\end{itemize}

\subsection{Estás bloqueado?}

Abre o cartão, escreve num comentário o que te está a impedir e menciona quem pode ajudar. Podes ainda destacá-lo com \btn{Sinalizar}, no menu de três pontos do cartão: fica realçado no quadro. Avisa também na daily. Pedir ajuda cedo não é falhar.

\subsection{Acabar um cartão}

Só passes um cartão a \textit{Feito} quando:

\begin{itemize}[label=\textcolor{azul}{$\square$}]
  \item cumpre os critérios de aceitação combinados com o Product Owner;
  \item o código foi revisto por outra pessoa da equipa (\textit{pull request} aprovado);
  \item tem testes e eles passam;
  \item está integrado na \texttt{main} sem quebrar nada;
  \item a documentação está atualizada e o Product Owner aceitou.
\end{itemize}

Se tens dúvidas de que está pronto, pergunta ao Product Owner antes de o arrastares.

\subsection{Criar um cartão novo}

\begin{enumerate}
  \item Carrega em \btn{Criar}, o botão azul no topo.
  \item Escolhe o tipo: \btn{Tarefa} para trabalho técnico ou de organização, \btn{História} para uma user story, \btn{Bug} para um erro.
  \item Escreve o \btn{Resumo}: curto, no formato «verbo + coisa». Por exemplo, \textit{Recuperar palavra-passe}.
  \item Na \btn{Descrição}, explica o que é para fazer. Se for uma história, usa «Como \dots, quero \dots, para \dots» e acrescenta os critérios de aceitação.
  \item Em \btn{Principal}, escolhe o epic a que pertence.
  \item Carrega em \btn{Criar}.
\end{enumerate}

Deixa o cartão sem sprint e sem responsável: o Product Owner trata da ordem no Backlog e, no Planning, decide-se quem faz o quê. Nunca apagamos cartões. Se foi engano, avisa o Scrum Master.

\begin{tcolorbox}[dica, title={Dica: caixas de verificação}]
  No editor da descrição, escreve \texttt{[]} seguido de um espaço para criar uma lista com caixas de seleção. É útil para os critérios de aceitação.
\end{tcolorbox}

\subsection{Estimar com story points}

Nos Plannings estimamos o esforço de cada cartão em \textit{story points}. Não são horas: medem quanto trabalho um cartão dá em comparação com os outros.

\begin{itemize}
  \item Usamos a escala \textbf{1, 2, 3, 5, 8 e 13}.
  \item Fazemos \textit{planning poker} simples: cada um pensa num número e todos o mostram ao mesmo tempo. Se houver muita diferença, quem escolheu o mais baixo e o mais alto explica porquê, e repete-se.
  \item Um cartão com mais de 13 pontos é grande demais. Parte-se em cartões mais pequenos.
\end{itemize}

Para registar, abre o cartão e preenche \btn{Story point estimate}.

\subsection{Escrever a daily}

Fazemos a daily como o contrato diz: 15 minutos, com o Quadro aberto. A seguir, \textbf{cada um escreve o seu registo} como comentário no cartão \textit{Daily – Sprint N} do sprint em curso (por exemplo o \textbf{EDS-4}, no Sprint 1). Cola este modelo:

\begin{codigo}
Dia 14/10 - Nome\\
Feito: o que fiz desde a última daily\\
A fazer: o que vou fazer hoje\\
Bloqueios: o que me está a impedir (ou 'nenhum')
\end{codigo}

Quem não pôde estar escreve na mesma, no mesmo dia. O Scrum Master confirma, até ao fim do dia, que está tudo escrito.

\subsection{Registar uma falta}

As faltas às aulas e às reuniões ficam no cartão \textit{Presenças – Sprint N} (o \textbf{EDS-19}, no Sprint 1). Cada falta é um comentário:

\begin{codigo}
Data - Aula ou reunião - Quem faltou - Avisou antes? (sim ou não) - Motivo
\end{codigo}

Quem falta também escreve o seu contributo no Jira nesse dia, como diz o contrato. No fim do sprint, o Scrum Master faz o resumo no mesmo cartão.

%--------------------------------------------------------------------
\section{Jira e GitHub juntos}

Cada cartão tem uma chave, por exemplo \texttt{EDS-13}. Se a escrevermos nos sítios certos, o Jira liga o trabalho ao cartão e a equipa percebe o que foi feito para cada coisa. A única branch do repositório é a \texttt{main}.

\begin{itemize}
  \item \textbf{Branch:} uma por cartão, a começar pela chave em maiúsculas: \texttt{feature/EDS-13-personas}.
  \item \textbf{Commit:} mensagem curta que diz o que mudou, com a chave no princípio: \texttt{EDS-13 acrescenta as personas}.
  \item \textbf{Pull request:} a chave vai no título e a descrição diz o que foi feito. Outra pessoa da equipa revê antes de integrar.
\end{itemize}

\begin{codigo}
git switch main\\
git pull\\
git switch -c feature/EDS-13-personas
\end{codigo}

A chave tem de estar em maiúsculas (\texttt{EDS-13}, e não \texttt{eds-13}). Passwords, chaves de API e dados pessoais nunca entram no repositório.

%--------------------------------------------------------------------
\section{Se és Product Owner ou Scrum Master}

Cada um de nós é Product Owner num sprint e Scrum Master noutro. A rotação está no contrato e na página \textit{Equipa e papéis} do Confluence. Isto é o que cada papel faz no Jira.

\subsection{Product Owner}

\begin{itemize}
  \item \textbf{Ordenar o Backlog.} No \btn{Backlog}, arrasta os cartões: o mais importante em cima. Só o Product Owner mexe na ordem; os outros sugerem e ele decide, depois de ouvir a equipa.
  \item \textbf{Escrever ou rever as user stories} e os critérios de aceitação.
  \item \textbf{Refinement}, a meio da semana: completar descrições, estimar e partir cartões grandes. Uma nota curta no cartão \textit{Planning} chega.
  \item \textbf{Aceitar ou recusar} os cartões na Review. O que não for aceite volta para \textit{A fazer} ou \textit{Fazendo}.
\end{itemize}

\subsection{Scrum Master}

\begin{enumerate}
  \item \textbf{Planning.} No \btn{Backlog}, confirma que o sprint existe (os cinco já estão criados). Com a equipa, arrasta para ele os cartões escolhidos e define o \btn{Responsável} de cada um.
  \item Carrega em \btn{Iniciar sprint}. Confirma o nome, ajusta as \textbf{datas} para as oficiais (o Jira preenche com o momento em que clicas) e confirma o objetivo do sprint, que já traz o PO e o SM.
  \item Preenche a ata no cartão \textit{Planning – Sprint N}: data, quem esteve, objetivo e cartões escolhidos. Junta uma captura do ecrã do sprint.
  \item \textbf{Durante o sprint}, confirma que as dailies estão escritas até ao fim de cada dia.
  \item \textbf{Review.} Preenche o cartão \textit{Review – Sprint N}: se o objetivo foi cumprido, o que o PO aceitou e o que volta ao backlog. Anexa capturas do Quadro e dos \btn{Relatórios}.
  \item \textbf{Retrospetiva.} Escreve no cartão \textit{Retrospetiva – Sprint N} o que correu bem, o que correu mal e duas ou três melhorias, cada uma com responsável.
  \item \textbf{Fechar.} Carrega em \btn{Concluir sprint}. O Jira pergunta o que fazer aos cartões que não ficaram \textit{Feito}: escolhe o \btn{Backlog}. Como o contrato diz, o que não ficou pronto volta ao Product Backlog e o sprint não se prolonga.
  \item Passa a \textit{Feito} os cartões de registo (Daily, Planning, Review, Retrospetiva e Presenças) que já estão preenchidos.
\end{enumerate}

\begin{tcolorbox}[atencao, title={Não inventamos registos}]
  Se um cartão de registo está vazio, não o passes a \textit{Feito}. Os docentes veem o histórico: o que escrevemos tem de ter acontecido mesmo, no dia em que aconteceu.
\end{tcolorbox}

\subsection{O que tem de ficar no Jira em cada sprint}

O enunciado avalia o Scrum através do que fica registado:

\begin{itemize}[label=\textcolor{azul}{$\square$}]
  \item Sprint com nome, datas oficiais e objetivo, com o PO e o SM indicados.
  \item Cartões do sprint com responsável, estimativa e descrição.
  \item Ata do Planning, com captura do ecrã.
  \item Um registo diário de cada pessoa, em cada dia útil.
  \item Nota do Refinement.
  \item Review, com captura do Quadro e dos relatórios.
  \item Retrospetiva, com as melhorias e os responsáveis.
  \item Presenças e faltas do sprint.
  \item Backlog atualizado, com o que ficou por fazer.
\end{itemize}

%--------------------------------------------------------------------
\section{Problemas comuns}

\tabelaini
\begin{tabularx}{\textwidth}{L{4.6cm} Y}
  \rowcolor{azul} \th{O que se passa} & \th{O que fazer} \\
  Não vejo o projeto & Ainda não aceitaste o convite, ou foste convidado com outro email. Diz ao Tiago o email que usas. \\
  Não consigo arrastar nem editar nada & Provavelmente tens a função de leitura. Pede ao Tiago para te passar a membro do projeto. \\
  O Quadro está vazio & O sprint em curso não tem cartões, ou estás a filtrar pelo teu nome e ainda não tens nenhum. Tira o filtro, ou vai ao Backlog. \\
  Um cartão desapareceu & Pode estar noutro sprint, noutro epic, ou escondido por um filtro. Usa a pesquisa (tecla \texttt{/}) com o nome ou a chave. \\
  Criei um cartão no sítio errado & Abre-o e corrige os campos \btn{Sprint} e \btn{Principal}, ou arrasta-o no Backlog. \\
  Apaguei um cartão sem querer & Não dá para desfazer. Por isso não apagamos cartões: se foi engano, avisa o Scrum Master. \\
  Recebo emails a mais & Nas definições pessoais da conta há opções de notificações. Podes reduzi-las, mas não ignores as menções com \texttt{@}. \\
  A ligação ao GitHub não aparece & Confirma se a chave está em maiúsculas (\texttt{EDS-13}) e se fizeste \textit{push}. \\
  Esqueci-me da palavra-passe & Usa a opção «Esqueceste-te da palavra-passe?» no ecrã de entrada da Atlassian. \\
  Não encontro um botão & Procura pelo que ele faz, não pelo nome. Os nomes mudam. Se continuar a faltar, pergunta no grupo. \\
  \bottomrule
\end{tabularx}
\tabelafim

%--------------------------------------------------------------------
\section*{Anexo: um painel com as secções do enunciado}

O enunciado fala em criar um «dashboard» com as secções \textit{Product Backlog}, \textit{To Do}, \textit{In Progress}, \textit{Done}, \textit{Artifacts}, \textit{Daily SCRUM} e \textit{Team Contract}. É opcional, mas dá para o ter num só ecrã: guarda um filtro por secção e põe cada um num painel.

\begin{enumerate}
  \item Abre a pesquisa avançada (\btn{Filtros}), cola a pesquisa da tabela, carrega em \btn{Guardar como} e dá o nome da secção.
  \item Depois abre \btn{Painéis}, cria um painel novo e acrescenta, para cada filtro, o gadget \btn{Resultados do filtro}.
\end{enumerate}

\tabelaini
\begin{tabularx}{\textwidth}{N{3.2cm} Y}
  \rowcolor{azul} \th{Secção} & \th{Pesquisa} \\
  Product Backlog & \small\texttt{project = EDS AND issuetype != Epic AND sprint is EMPTY AND statusCategory != Done ORDER BY Rank} \\
  To Do & \small\texttt{project = EDS AND sprint in openSprints() AND statusCategory = 'To Do'} \\
  In Progress & \small\texttt{project = EDS AND sprint in openSprints() AND statusCategory = 'In Progress'} \\
  Done & \small\texttt{project = EDS AND sprint in openSprints() AND statusCategory = Done} \\
  Artifacts & \small\texttt{project = EDS AND parent = EDS-2} \\
  Daily SCRUM & \small\texttt{project = EDS AND parent = EDS-3} \\
  Team Contract & \small\texttt{project = EDS AND parent = EDS-1} \\
  \bottomrule
\end{tabularx}
\tabelafim

\end{document}
